\section*{Thème de la mission : objectifs et contexte}
\addcontentsline{toc}{section}{Thème de la mission : objectifs et contexte}

Ce rapport de Mission 4 porte sur l'analyse de la politique de \glos{RSE} (Responsabilité Sociétale des Entreprises) et sur le positionnement du \textit{numérique responsable}\footnote{On entend par numérique responsable l'ensemble des pratiques visant à réduire l'empreinte environnementale et sociale des systèmes d'information tout en garantissant leur conformité éthique.} d'Iveco Group. Placé au croisement de la stratégie groupe et des réalités opérationnelles du site de Haute-Goulaine, ce travail s'appuie exclusivement sur les documents fournis par l'entreprise et accessibles dans le cadre de la mission.

\subsection*{Objectifs}

Les objectifs assignés à cette mission sont les suivants :

\begin{itemize}
    \item Identifier et restituer de manière rigoureuse les engagements, les organes de gouvernance et les indicateurs clés (\glos{KPI}\footnote{Key Performance Indicator : indicateur clé de performance.}) de la politique RSE d'Iveco Group.
    \item Analyser la mise en \oe{}uvre sociale, en particulier les politiques de Diversité, Équité et Inclusion (\glos{DEI}\footnote{Diversity, Equity and Inclusion : diversité, équité et inclusion.}), de santé et de sécurité, ainsi que les initiatives communautaires.
    \item Faire le lien entre ces engagements groupe et l'environnement informatique du site de Haute-Goulaine, en proposant des recommandations concrètes, chiffrables et réalistes pour l'équipe IT.
\end{itemize}

\subsection*{Périmètre et méthodologie}

L'analyse concerne Iveco Group dans sa globalité, en particulier les données publiées pour l'exercice 2025 (rapport annuel, \textit{GHG Statement}\footnote{Greenhouse Gas (gaz à effet de serre) Statement (déclaration).}, documents thématiques RSE/DEI et politique de confidentialité). Le site de Haute-Goulaine sert de terrain d'observation concret pour illustrer les leviers d'action du numérique responsable.

La méthode adoptée repose sur trois étapes :

\begin{enumerate}
    \item \textbf{Extraction documentaire~:} relevé des faits, chiffres et engagements dans les ressources mises à disposition.
    \item \textbf{Analyse critique~:} mise en perspective des indicateurs, identification des points forts et des axes de progrès.
    \item \textbf{Recommandations opérationnelles~:} traduction des engagements groupe en actions adaptées à l'équipe IT locale.
\end{enumerate}

\subsection*{Plan du rapport}

Le développement s'articule autour de trois parties, avant une conclusion qui en tire le bilan et les perspectives :

\begin{enumerate}
    \item \textbf{Gouvernance et engagements globaux~:} stratégie de durabilité, gouvernance, labels et dispositif conformité.
    \item \textbf{La mise en \oe{}uvre sociale~:} DEI, santé-sécurité et engagement communautaire.
    \item \textbf{Numérique responsable et recommandations~:} diagnostic du lien RSE/IT et propositions chiffrées pour l'équipe de Haute-Goulaine.
\end{enumerate}
